% ==========================================================================
% GENERATED by scripts/split_paper.py. Do not edit this file.
%
% Source: imported_papers/Neural_Computation__CNA/NECO-D-24-00304-Source.tex
% Part 11: Normalization
%
% Edits belong upstream, in the published paper. To change what the
% chapter includes, change imports/cna.toml and re-run
% the splitter.
% ==========================================================================
\section{Normalization}\label{sec:normalization}
The pre-synaptic activity in $S2$ is normalized using a two-step approach that confines the neuronal activations $\boldsymbol{a}^{\mathcal{I}(S2)}$ to the interval $(0,1)$. 

We first apply a size limit:
\begin{equation}
  \boldsymbol{a}^{\mathcal{I}(S2)}_{c,t} = \begin{cases}
   		\boldsymbol{a}^{\mathcal{I}(S2)}_{c,t}, & \text{ if } \boldsymbol{a}^{\mathcal{I}(S2)}_{c,t} \leq \lambda \\
   		\lambda - \frac{1}{2}(\boldsymbol{a}^{\mathcal{I}(S2)}_{c,t} - \lambda), & \text{otherwise}.
  	\end{cases}\label{eq:saturation}
\end{equation}
According to this formula, $\boldsymbol{a}^\mathcal{I}$ grows undiminished with $\boldsymbol{a}^\mathcal{I}$ until it reaches the value $\lambda$, and thereafter actually diminishes with slope $-1/2$. (In our experiments we found the value $\lambda = 1.3 \cdot \frac{h + w}{2}$ to work well). 
This saturation helps to mitigate imbalances between net fragments with differing numbers of participating neurons. In the next step, we introduce the normalization
\begin{equation}
    \boldsymbol{a}^{\text{norm }(S2)}_{c,i,t} = \max \left( \frac{\boldsymbol{a}^{\mathcal{I}(S2)}_{c,i,t}}{\max_{i'} \boldsymbol{a}^{\mathcal{I}(S2)}_{c,i',t}}
    \label{eq:normalization}, 0 \right)
\end{equation}
Here, $i$ refers to spatial locations within the image, i.e., $i=(h,w)$ where $h \in H, w \in W$.
This normalization confines activations to the interval $[0,1]$, where neurons receiving maximal support are mapped to $1$, while those with less support are mapped to correspondingly smaller values. 
This normalization is relative to the maximal value that itself is growing during learning as it is important to give neurons initially (when recurrent connections are still small) a chance to fire and later when recurrent weights between correlated patterns have been learned, to suppress those neurons that are part of noise patterns. 
